\newsection{Mengen, Relationen und Abbildungen}{task-1}

\begin{enumerate}
  \item {
    Es seien N, M zwei Mengen mit $|N| = n, |M| = m$ und $|N \setminus M| = d$.
    Geben Sie die folgenden Kardinalitäten an.
    Begründen Sie Ihre Antwort kurz, indem Sie die Zwischenschenschritte angeben.
    \begin{enumerate}
      \item $|\mathcal{P}(M \cup N)| = $ \vspace*{4cm}
      \item $|(M \cap N) \times M| = $ \vspace*{4cm}
    \end{enumerate}
  }
  \item {
    Zeigen Sie mittels vollständiger Induktion, dass für alle $n \in \mathbb{N}$
    gilt:
    \[
      \sum_{i=0}^{n-1} (3i^2 + 3i + 1) = n^3
    \]
  }
  \newpage
  \item {
    Kreuzen Sie für jede der folgenden Funktionen an, ob sie \textit{injektiv},
    \textit{surjektiv} oder \textit{weder noch} ist.
    Eine Begründung ist \textit{nicht} notwendig.

    \setlength{\tabcolsep}{10pt}
    \renewcommand{\arraystretch}{2}
    \begin{tabular}{@{}l c c c@{}}
      & $f$ injektiv & $f$ surjektiv & $f$ weder injektiv noch surjektiv \\
      $f: \mathbb{R} \to \mathbb{R} $ mit $f(x) = x^2 - 2x$ & $\square$ & $\square$ & $\square$ \\
      $f: \mathcal{P}(\mathbb{N}) \to \mathbb{N}_0 $ mit $f(A) = |A|$ & $\square$ & $\square$ & $\square$ \\[10pt]
      $f: \mathbb{R} \to \mathbb{R} $ mit $f(x) =
        \begin{cases}
          -(x^2), & x \leq 0 \\
          x^2 & x > 0
        \end{cases}
        $ & $\square$ & $\square$ & $\square$ \\[10pt]
      $f: \mathbb{R}^2 \to \mathbb{R}^2 $ mit $f(x, y) = (y-7, \frac{1}{4} x)$ & $\square$ & $\square$ & $\square$ \\
      $f: \mathbb{R}^2_{\geq 0} \to \mathbb{R}_{\geq 0} $ mit $f(x, y) = x^y$ & $\square$ & $\square$ & $\square$ \\
      $f: \mathbb{N}^2 \to \mathbb{N} $ mit $f(x, y) = x^2 + y^2$ & $\square$ & $\square$ & $\square$
    \end{tabular}
  }
  \vspace{\fill}
  \item {
    Seien $M = \{A, B, C, D\}$ eine Menge und $R \subseteq M \times M$ die unten
    schematisch dargestellte Relation auf $M$. Geben Sie jeweils mit
    \textit{kurzer} Begründung an, ob $R$ reflexiv, symmetrisch, transitiv,
    antisymmetrisch bzw. eine Äquivalenzrelation ist.

    \begin{center}
    \begin{tikzpicture}[node distance=2cm, auto]
      \node[state] (A)  {$A$};
      \node[state] (B) [right=of A]  {$B$};
      \node[state] (C) [below=of A]  {$C$};
      \node[state] (D) [below=of B]  {$D$};
      \path[->]
      (A) edge [bend left =15] node {} (B)
      (B) edge [bend right =-15] node {} (A)
      (A) edge [bend left =15] node {} (C)
      (C) edge [bend right =-15] node {} (A)
      (B) edge [bend left =15] node {} (D)
      (D) edge [bend right =-15] node {} (B)
      (B) edge [bend left =15] node {} (C)
      (C) edge [bend right =-15] node {} (B);
    \end{tikzpicture}
    \end{center}
    \vspace{\fill}

    \begin{minipage}{\linewidth}
      reflexiv: \hrulefill \\[10pt]
      symmetrisch: \hrulefill \\[10pt] \hrule \vspace{18pt}
      transitiv: \hrulefill \\[10pt] \hrule \par\vspace{18pt}
      antisymmetrisch: \hrulefill \\[10pt] \hrule \par\vspace{18pt}
      Äquivalenzrelation: \hrulefill
    \end{minipage}
  }
  \vspace{\fill}
\end{enumerate}